\par\bigskip\noindent\textbf{K99 supporting proof: An independent evaluation of the compact surface intersections}\par\nobreak
We use the actual three-end atlas and the compact surface components
recorded in \path{code/part3/e6_order_three_fixed_surfaces.json.} This audit
computes their normal intersections by a second exact procedure. The
four lower surfaces at p=1 and the exterior reference doubles remain
outside its scope.

In common characters the auxiliary circle has cocharacter
\(\epsilon=(0,1,0)\): it scales \(Z\) and fixes
\(V_{\rm reg},P,p\). The actual native cocharacters in the three
ends are
\[
 \epsilon_0=(-3,4,0,0),\qquad
 \epsilon_1=(1,1,0,0),\qquad
 \epsilon_\infty=(0,1,0,0).
\]
They follow by applying the inverse native-to-common matrices. The
period matrices fix these cocharacters. Their complete coefficient
functions depend only on the unchanged base variables, so the full
period maps commute with this action. The monomial chart morphisms
therefore extend it to every chart of these compact surfaces and to
their actual physical normal lines. We use only this explicitly
extended action; no auxiliary action on the external attachment is
assumed.

Each compact surface has two auxiliary-fixed sections, both actual
copies of \(\mathbb P^1_p\). For a section \(S\), let \(k\)
be the nonzero auxiliary weight on its normal line in the surface,
let \(s=S^2\) be that line\textquotesingle s degree, and let \(m_i,d_i\) be
the auxiliary weight and the degree of the i-th physical normal
line on \(S\). All these integers are read from the original
free ray matrices and the three full lower-ray clutches. The two
values of \(k\) on each surface are opposite. The 113 ruled
surfaces have no other auxiliary-fixed components. Each of the four
surfaces with an additional infinity chart has one further isolated
auxiliary-fixed point. Its two tangent weights and its physical
normal-line weights are read directly from that infinity chart.
The program tests these assertions for all 117 component stars.

Write \(x_i\) for the ordinary Chern class of the i-th line on
the section, and \(y\) for the Chern class of its normal line in
the surface. Expanding the actual localization denominator gives
\[
 \frac{(x_i+m_i u)^2}{y+ku}
 =\frac{m_i^2u}{k}+\frac{2m_ix_i}{k}
       -\frac{m_i^2y}{k^2}+O(u^{-1}).
\]
Integration on \(S\) therefore gives the contribution
\(2m_id_i/k-m_i^2s/k^2\) to the line square. Likewise the
mixed product has contribution
\[
 \frac{m_1d_2+m_2d_1}{k}-\frac{m_1m_2s}{k^2}.
\]
At an isolated auxiliary point, the corresponding terms are
\(m_i^2/(k_1k_2)\) and \(m_1m_2/(k_1k_2)\), using its two
actual tangent weights. Thus this calculation uses the literal
section/point representations and the full line clutches; it does
not invert the surface intersection matrices used by the first
certificate.

The tangent calculation is independent as well. Along a fixed
section, the surface tangent bundle has roots \(c_1(TS)\) and
\(y+ku\). The degree-four term \(c_1(TS)^2\) vanishes on the
curve. Dividing the remaining square by \(y+ku\) and integrating
gives \(s\). At an isolated point the Pontryagin contribution is
\((k_1^2+k_2^2)/(k_1k_2)\). The two section self-intersections
sum to zero for each ruled surface and to \(-1\) for each
one-point blowup. In the latter cases the actual isolated point has
opposite tangent weights, and hence contributes \(-2\). The
resulting Pontryagin numbers are zero and \(-3\), respectively,
with the original complex orientations.

The executable audit recomputes all three normal intersections and
the tangent Pontryagin number for every surface. Each normal
intersection is an integer and agrees with the first integral
intersection-matrix computation. Substitution in the complete SZ
surface polynomial gives the same physical inner pair
\[
 (k_M,\ell_M)=\left(\frac{1121}{18},\frac{596}{3}\right),
 \qquad \frac{4k_M-\ell_M}{24}=\frac{227}{108}.
\]
The certificate passes 1,644 exact assertions and retains every
section weight, line degree, section self-intersection and extra
isolated-point representation. Its input hashes identify all three
complete fan certificates used.

Adding the independently constructed 276 compact curves and the
two ten-point corner blocks gives the specified compact subtotal
\[
 \left(-\frac{1895}{18},-\frac{50}{3}\right)
 +\left(\frac{1121}{18},\frac{596}{3}\right)
 +2(-1,-31)=(-45,120),
\]
whose displayed inner residue is \(-25/2\). The p=1 lower
surfaces, exterior reference doubles, and outer-boundary generator
and orientation maps are not included in this equality. Their
calculation is required before assigning the characteristic return.

